\documentclass[11pt,letterpaper]{article}
\usepackage[margin=0.8in]{geometry}
\usepackage[T1]{fontenc}
\usepackage{lmodern,microtype,amsmath,booktabs}
\usepackage[hidelinks]{hyperref}
\hypersetup{pdfcreator={},pdfproducer={}}
\pdfinfoomitdate=1
\pdfsuppressptexinfo=15
\setlength{\parindent}{0pt}
\setlength{\parskip}{6pt}
\begin{document}
\begin{center}
{\large\bfseries Supporting Information for\\[5pt]
Identifying Solar and Heliospheric Fronts With Conservation Laws and the MHD Riemann Problem}\par
\vspace{5pt}
Olena Podladchikova
\end{center}

\textbf{Contents:} Movies S1--S4 and an accompanying animation package
containing the rendering script, retained inputs, and provenance records.
The movies accompany the solar-image panels in Figure~1 and the SUVI and
AIA comparisons in Figures~6 and~7 of the article.

All movies show observed images with fixed display scales within each
sequence. The left panel shows brightness or radiance; the right panel
shows a difference image. Red and blue indicate positive and negative
changes. Playback is accelerated, and the printed UT timestamps give
the observation times. Each exposure is held on screen; no image is
interpolated between exposures. Previously calculated feature tracks
and model positions are retained without refitting.

\textbf{Movie S1. The 2010 June 13 AIA example (E11).}

Sixteen AIA 193~\AA\ exposures from 05:38:08.07 to 05:41:08.07~UT
show the evolving off-limb structure, at approximately 12~s spacing.
The right panel subtracts the reference at 05:35:08.07~UT.
The retained displays use spatial Gaussian smoothing with a standard
deviation of 1.20~arcsec and a difference scale of $\pm12$~DN~s$^{-1}$.
No temporal smoothing is applied.
Turquoise and orange curves mark the retained inner and outer
brightness-crest fits. A curve appears only in exposures for which
that fit exists; it is not extrapolated into other frames.
The curves describe image features. The conditional fast-shock test
in the article concerns a prescribed plasma-state pair motivated by
the observations discussed by Kozarev et al.\ (2011) and Ma et al.\ (2011).

\textbf{Movie S2. The 2011 February 16 AIA example (E05).}

Twenty-nine AIA 193~\AA\ exposures from approximately 14:25:32.84
to 14:31:08.84~UT show the region used for the EIS comparison.
The images are exposure-normalized and registered on a common grid.
The first selected exposure supplies the base-difference reference;
its difference panel is therefore zero. Yellow boxes show the sixteen
retained alternatives for placement of the EIS sampling region.
They do not trace one uniquely resolved front. The timestamps use
the retained time offsets, anchored to the article exposure at
14:28:44.84~UT, and are rounded for display. The observational context
is the analysis of Veronig et al.\ (2011). The two conditional state
pairs in the article explore different interpretations of the
incomplete plasma constraints.

\newpage
\textbf{Movie S3. The SUVI geometric test on 2017 September 10.}

Five SUVI 195~\AA\ exposures run from 15:54:24.75 to 15:58:24.76~UT.
Their actual spacings are approximately 70, 50, 50, and 70~s.
The left panel shows calibrated radiance on a logarithmic scale.
The right panel subtracts the 15:30:24.69~UT reference after sampling
on the retained common coordinate grid. It uses a fixed asinh color
scale. No additional spatial or temporal smoothing is applied here.

Orange C marks the selected outer bright feature and pink E the
selected inner edge. The dotted yellow outline marks the fixed
trailing aperture B. At each time, the tested association places the
P100 fast shock at C and its contact at E. One common shift and scale
then determine the purple fast-rarefaction interval. This interval
lies outside B from the second exposure onward. The labels show the
model association being tested; they do not establish those wave
types from image brightness. The movie visualizes the same geometric
construction as the article. SUVI observations of this eruption were
presented by Seaton et al.\ (2018).

\textbf{Movie S4. The AIA comparison along G1 on 2017 September 10.}

Five AIA 193~\AA\ exposures run from 16:00:41.84 to 16:05:29.84~UT,
at approximately 72~s spacing. The right panel subtracts each
preceding exposure; the first uses 15:59:29.84~UT as its reference.
These images are retained display levels used to compare positions
and shapes. They are not used to infer density or temperature.

The gray curve is G1, an assumed great-circle path on a sphere of
radius 696~Mm, projected into the image plane. The open turquoise
circle marks the selected image maximum. Colored marks show the P100
wave positions: orange for the fast shock, pink for the contact,
purple for the fast rarefaction, and green for the other model waves.
Only the first fast-shock position is aligned with the selected
maximum. All later positions follow from the fixed model speeds,
the assumed initial model age of 120~s, and the same geometry.
The movie shows outward feature motion together with the positions
at which the other model waves can be sought. The contact and
rarefaction remain predictions requiring independent plasma tests.

\textbf{Reproduction and data credits.}

The code archive provides \texttt{run.py},
\texttt{requirements.txt}, and all retained movie inputs.
After installing the dependencies, run
\texttt{python run.py animations}. Python and Matplotlib generate
the annotated figures; FFmpeg encodes MP4 files. GIF copies and
representative PNG/PDF frames are also generated. The README explains
the inputs and command; \texttt{MANIFEST.json} records file checksums.
The rendering code and numerical calculations are included in the
release archived at
\url{https://doi.org/10.5281/zenodo.23084162}.

AIA data are courtesy of NASA/SDO and the AIA science team.
SUVI data are provided by NOAA/NCEI. References cited in these captions
are listed in the article bibliography.
\end{document}
